\documentclass[a4paper,12pt]{article}

\usepackage[french]{babel}
\usepackage[T1]{fontenc} 
\usepackage[latin1]{inputenc}
%\usepackage{amsmath}

\author{R�mi Synave, Stefka Gueorguieva, Pascal Desbarats}
\title{Manuel de la biblioth�que VECTOR}
\date{}

\begin{document}
\maketitle

\section{Utilisation}
Cette biblioth�que impl�mente une structure de vecteur 3D les fonctions associ�s pour les manipuler ainsi que des op�rations comme la normalisation, le produit scalaire ou vectoriel.\\

\section{Structures de donn�es}

Le vecteur 3D est une structure comportant trois r��ls : la composante suivant X, la composante suivant Y et la composante suivant Z.\\
\begin{verbatim}
typedef struct
{
  double dx,dy,dz;
}vector3d;
\end{verbatim}

\section{Fonctions}


\textbullet void vector3d\_init(vector3d *v, double dx, double dy, double dz)\\
Initialisation d'un vecteur 3D avec trois r��ls.\\


\textbullet void vector3d\_display(vector3d v)\\
Affichage d'un vecteur 3D.\\


\textbullet double vector3d\_size(vector3d v)\\
Calcul de la taille d'un vecteur 3D.\\


\textbullet int vector3d\_equal(vector3d v1, vector3d v2)\\
Retourne 1 si les composantes des deux vecteurs sont �gales une � une.\\


\textbullet void vector3d\_normalize(vector3d *v)\\
Normalisation du vecteur 3D. Chaque composante est divis�e par la taille du vecteur 3D.\\


\textbullet double vector3d\_scalarproduct(vector3d v1, vector3d v2)\\
Calcul du produit scalaire de deux vecteurs 3D.\\


\textbullet vector3d vector3d\_vectorialproduct(vector3d v1, vector3d v2)\\
Calcul du produit vectoriel de deux vecteurs 3D.\\


\textbullet void vector3d\_reverse(vector3d *v)\\
Inversion du vecteur 3D. Chaque composante est multipli� par -1.\\

\end{document}
